\subsection*{Cap\'itulo 4: C\'alculo}
\textbf{4.1}\label{resp:4-1}\quad
(1) 6. (2) 2. (3) 3. (4) Por la derecha $+\infty$, por la izquierda $-\infty$: no existe. (5) $\frac{(x+2)(x-1)}{x-1}=x+2\to3$. (6) $\frac52$. (7) 8; la tabla da $7{,}9$; $7{,}99$; $7{,}999$ por la izquierda y $8{,}001$; $8{,}01$; $8{,}1$ por la derecha. (8) $\frac{1}{\sqrt{x+4}+2}\to\frac14$.

\medskip\noindent\textbf{4.2}\label{resp:4-2}\quad
(1) $3$. (2) $3x^2-2$ en $x=2$ da $10$. (3) $f'(x)=e^x(1+x)$ y $f'(0)=1$. (4) $6(2x+1)^2$ en $x=1$ da $54$. (5) $y=3x-2$. (6) $\frac{x^2-2x-1}{(x-1)^2}$. (7) $v(3)=32$ m/s y $a=10$ m/s$^2$. (8) $\frac{2x}{x^2+1}$.

\medskip\noindent\textbf{4.3}\label{resp:4-3}\quad
(1) $f_x=2xy$, $f_y=x^2+3$. (2) $ye^{xy}$. (3) $\nabla f=(6, 8)$ y $\|\nabla f\|=10$. (4) $(0, 0)$ y $(1, 1)$; en $(0, 0)$ es silla ($D=-9$) y en $(1, 1)$ es m\'inimo ($D=27$, $f_{xx}=6$). (5) M\'inimo en $(2,-3)$ con valor $-13$. (6) Punto silla. (7) $-12$ y $-6$. (8) En la direcci\'on de $\nabla f=(y, x)=(3, 2)$.

\medskip\noindent\textbf{4.4}\label{resp:4-4}\quad
(1) $x=3$, m\'inimo con $f=-4$. (2) $x=2$, m\'aximo con $f=4$. (3) $x=4$, $C=4$. (4) $x=-1$ m\'aximo ($f=2$) y $x=1$ m\'inimo ($f=-2$). (5) Cuadrado de $15\times15$, \'area $225$ m$^2$. (6) $n=10$, $C=60$. (7) $x=10$, $G=100$. (8) Porque en cada paso se mueve contra el gradiente, que es la direcci\'on de mayor descenso, hasta llegar a un punto donde $\nabla J=0$.

\medskip\noindent\textbf{4.5}\label{resp:4-5}\quad
(1) $x^3+C$. (2) $2$. (3) $x^2+x+C$. (4) $12$. (5) $e-1$. (6) $2$. (7) $60$. (8) Riemann: $3{,}75$; exacto: $\frac83\approx2{,}667$.

\medskip\noindent\textbf{Mini proyecto:} (1) $8$ ms; sustituir da $0/0$ y hay que factorizar. (2) $x=4$ con $C(4)=4$ y $C''=2>0$. (3) $N(5)=60$.

\medskip\noindent\textbf{Autoevaluaci\'on 4:}\label{resp:q4}\quad 1 c, 2 b, 3 b, 4 b, 5 b, 6 b, 7 b, 8 b, 9 a, 10 c.
